% Ad Familiares 10.7 (ad-familiares-10-07)
% Source: https://raw.githubusercontent.com/PerseusDL/canonical-latinLit/master/data/phi0474/phi056/phi0474.phi056.perseus-lat1.xml
% Fetched by scripts/fetch_latin.py

% Dateline (Perseus): Scr. in Gallia Transalpina paulo post med. m. Mart. a. 711 (43) .

\ciceroLetterOpener{PLANCVS CICERONI}

\ciceroSection{1} plura tibi de meis consiliis scriberem rationemque omnium rerum redderem verbosius, quo magis iudicares omnia me rei p. praestitisse, quae et tua exhortatione excepi et mea adfirmatione tibi recepi (non minus enim a te probari quam diligi semper volui, nec te magis in culpa defensorem mihi paravi quam praedicatorem meritorum meorum esse volui) ; sed breviorem me duae res faciunt, una, quod publicis litteris omnia sum persecutus, altera, quod M. Varisidium, equitem R., familiarem meum, ipsum ad te transire iussi ex quo omnia cognoscere posses.

\ciceroSection{2} non medius fidius mediocri dolore adficiebar cum alii occupare possessionem laudis viderentur sed usque mihi temperavi dum perducerem eo rem, ut, dignum aliquid et, consulatu meo et vestra : exspectatione efficerem. quod spero, si me fortuna non fefellerit, me consecuturum ut, maximo praesidio rei p. nos fuisse et, nunc sentiant, homines et, in posterum memoria teneant. A te peto ut, dignitati meae suffrageris et, quarum rerum spe ad laudem me vocasti, harum fructu in reliquum facias alacriorem. non minus posse te quam velle exploratum mihi est. fac valeas meque mutuo diligas.
